\chapter{Groups 1 and 2: Elements That Readily Give Up Electrons}

\begin{kaichang}
Search the medicine cabinet and you will probably find calcium tablets. Their ingredients may include calcium carbonate and a little vitamin D. They sit quietly in the bottle, and you can chew and swallow them without anything happening, exactly as intended.

Now imagine a different scene, seen only in an online video or teacher demonstration: \textbf{never try this yourself}. A small piece of sodium is lifted from kerosene with forceps, blotted with filter paper, and dropped into water. It immediately melts into a bright silver bead, hissing and darting across the surface, sometimes bursting into yellow flame with a bang or even shattering the glass.

Sodium and calcium are neighboring elements that readily surrender electrons. Why must one be imprisoned in oil while the other is mild enough for a medicine bottle? Is swallowed calcium the same as calcium metal?

Make a guess. We will return to it at the end.
\end{kaichang}

\section{Three Brothers on Video}

First watch a public demonstration video of alkali metals reacting with water, choosing reputable science channels or school lessons. Three half-filled beakers receive rice-grain-sized pieces of lithium, sodium, and potassium. You will see:

\begin{itemize}[itemsep=2pt]
\item \textbf{Lithium} floats, moves slowly, and hisses with bubbles, like melting ice;
\item \textbf{Sodium} reacts much faster, its released heat melting it into a bead (its melting point is only $97.8\,^{\circ}\mathrm{C}$) that darts about with sharp bubbling sounds;
\item \textbf{Potassium} ignites in a pale violet flame upon contact, often with small popping sounds.
\end{itemize}

Look farther down the family: rubidium and cesium explode violently on touching water, so even formal demonstrations increasingly avoid them. Both are extremely expensive and dangerous and seldom encountered in everyday life.

What do the metals look like beforehand? Submerged in kerosene or paraffin oil like canned fruit, their surfaces look dull gray when removed. A knife cut, easy because they are as soft as firm cheese, reveals silvery metal. Its shine lasts only seconds before fading, like a cut apple discoloring in air.

Why oil? Air is a reactive soup for them. Oxygen takes electrons, making surface oxides and dulling the cut; water vapor is still more troublesome. Containment needs a liquid without oxygen or water that does not react and excludes air. Kerosene fits, with another convenient property: sodium's density is about \ud{0.97}{g/cm^{3}}, kerosene's around \ud{0.8}{g/cm^{3}}, so sodium sinks completely beneath the oil. Lithium is an exception: at only \ud{0.53}{g/cm^{3}}, it floats in kerosene and touches air, so it must be enclosed in solid paraffin. Even storage answers lie in physical properties.

Group 2's beryllium, magnesium, calcium, strontium, and barium, the alkaline-earth metals, are milder but still active characters. In a teacher's magnesium-burning demonstration, a short gray-silver ribbon held with crucible tongs ignites with white light too dazzling to stare at, like lightning entering the classroom. After tens of seconds, only light white powder remains, easily blown away. This was the light source of early photographic flashes and illumination flares. Calcium metal in water, less violent than sodium, steadily emits bubbles and gradually clouds the water. The sparingly soluble calcium hydroxide forms a faint milk of lime.

Meanwhile, eggshells, chalk, marble windowsills, and calcium tablets containing Group 2 elements sit quietly, even immersed in water.

That is our mystery: \textbf{why are the same two groups' elements sometimes violent and sometimes mild? Where exactly does the difference lie?}

\section{Valence Electrons: The Atom's Outermost Hand}

Chapter 9 described electron floors; Chapter 10 connected outer-electron counts to personality. This family's character grows from that rule.

Lithium, sodium, and potassium each have \textbf{1} outer electron; beryllium, magnesium, and calcium each have \textbf{2}. These one or two occupy the highest, outermost shell, farthest from the nucleus and screened by inner electrons (Chapter 11's shielding), making them the atom's most loosely held electrons.

Now bring in water. Water molecules are formidable partners, pulling at surface electrons (Chapter 23 explains their power). Sodium's account runs as follows: losing its outer electron costs some energy ransom, but surrounding the resulting positive ion with water molecules releases more, leaving a net gain. Hydrogen, another product, also carries heat away. If total energy falls, change may occur. The atom has no desire; it simply rolls downhill in energy.

What does electron loss produce? An \textbf{ion}, a charged atom. Sodium loses one electron and gains one positive charge; magnesium loses two and gains two. These ions have full outer floors, low in energy and stable. That is why they lose one or two electrons rather than three.

\begin{qiaoliang}
What proves only one is lost? \textbf{Ionization energies}, the energy needed to remove gaseous atomic electrons one by one. For sodium, per mole of atoms:
\begin{itemize}[itemsep=1pt]
\item First electron: about \ud{496}{kJ/mol};
\item Second: about \ud{4562}{kJ/mol}, over nine times dearer;
\item Third: about \ud{6912}{kJ/mol}, still expensive.
\end{itemize}
The first is cheap; the second jumps almost an order of magnitude because removal now attacks the full second shell, closer and more tightly bound. Magnesium tells the same story: about \ud{738}{} and \ud{1451}{kJ/mol} for the first two, then about \ud{7733}{kJ/mol} for the third. The price cliff precisely identifies the preferred electron loss: one for Group 1, two for Group 2, each stopping at its cliff. Thus common ion charges can be \textbf{read directly from group number}: Group 1 $+1$, Group 2 $+2$.
\end{qiaoliang}

\section{Why Each Brother Is More Hot-Tempered}

Lithium circles, sodium melts into a bead, potassium catches fire. Why does one family personality grow fiercer down the generations?

Take Chapter 11's conclusion: downward in a group, outer electrons occupy higher floors farther from the nucleus behind thicker inner-electron blankets, weakening attraction. Looser holding means quicker loss. Lithium's electron takes more effort to pull away; potassium's nearly falls off on contact. The water reaction's starting threshold falls and its violence rises.

Do not glide past a distinction: \textbf{readiness to lose electrons describes tendency, or reactivity; reaction violence describes speed}. They are different. Energy accounting determines tendency (Chapters 27 and 28), while conditions determine rate (Chapter 29). Three conditions particularly matter:

\begin{itemize}[itemsep=2pt]
\item \textbf{Temperature}: magnesium barely reacts with cold water at room temperature but bubbles steadily in hot water; ignition lets ribbon burn in air. Higher temperature energizes collisions enough to cross the reaction barrier.
\item \textbf{Surface area}: a whole calcium piece merely bubbles in water, but powder may react violently. Powdered metal has a front line nearly everywhere.
\item \textbf{The surface shell}: many active metals carry a thin oxide raincoat keeping water out. Sodium's freshly cut surface darkens as it sews this coat. Aluminum, Group 13 and revisited in Chapter 21, is more active than magnesium yet makes drink cans and aircraft because of its dense, strong oxide film.
\end{itemize}

Active therefore does not mean exploding at every encounter. A smooth aluminum nail in water does nothing, not because electron surrender is unfavorable but because the raincoat blocks access. Remember this for corrosion and batteries later.

\section{Now for the Symbols}

Observations, particles, and rules are in place. Chemical equations now compress their accounts into a line.

Sodium with water, general for Group 1 by replacing \chem{Na} with \chem{Li} or \chem{K}:
\[\chem{2Na} + \chem{2H_2O} \rightarrow \chem{2NaOH} + \chem{H_2}\uparrow\]
Read it as two sodium atoms each transferring one electron to water and becoming sodium ions. Water yields hydrogen gas, while the remainder combines with sodium ions as dissolved sodium hydroxide, a strong base discussed in Chapter 32. The demonstration's resulting water is hot and strongly alkaline to an indicator, with a soapy slippery feel (\textbf{do not test it by hand}; sodium hydroxide corrodes skin). How is the gas identified? A classic classroom demonstration collects bubbles, or makes soap bubbles through a delivery tube dipped in soap solution, and touches them with a burning splint. A small pop is hydrogen mixed with air burning, the squeaky-pop test. The evidence chain closes: the gas is hydrogen; the equation's \chem{H_2} is not a guess.

Magnesium burning in air:
\[\chem{2Mg} + \chem{O_2} \rightarrow \chem{2MgO}\]
The dazzling light comes from energy released when magnesium and oxygen form magnesium oxide, bond-formation energy discussed in Chapters 17 and 27. Fireworks and flares add magnesium powder for this white light. Strontium salts give red, calcium orange-red: Group 2 shows family ties in flame colors too.

Calcium with water, Group 2's gentler version:
\[\chem{Ca} + \chem{2H_2O} \rightarrow \chem{Ca(OH)_2} + \chem{H_2}\uparrow\]

Estimate the energy scale. Sodium reacting with water releases about \ud{180}{kJ} per mole. A small demonstration piece of about \ud{2.3}{g}, 0.1 mole, releases roughly \ud{18}{kJ}. If all heats \ud{50}{g} of water, with \ud{4.2}{J} needed per gram per $1\,^{\circ}\mathrm{C}$, the temperature rises about $85\,^{\circ}\mathrm{C}$. Rice-sized metal pushes half a cup from room temperature toward boiling within seconds, incidentally igniting hydrogen. Now potassium's immediate flame makes sense.

\section{Davy's Electric Shovel Unearths Two New Elements}

\begin{zhengju}
Early nineteenth-century chemistry listed irreducible substances, then considered elements. Soda, sodium carbonate, and potash, potassium carbonate, appeared prominently. But Lavoisier had wondered whether their sibling-like personalities meant they were compounds of unknown metals rather than basic substances.

Suspicion was not enough: nobody could break them apart. The hottest fires and strongest acids left soda and potash unchanged. In 1807, Humphry Davy of Britain's Royal Society obtained a new weapon, the world's largest voltaic battery, hundreds of copper--zinc pairs forming an electric shovel.

His idea was direct: use electricity to tear apart what chemistry could not. Current supplies electron-moving force; if these substances join metal and other components by electron transfer, a stronger electron flood might separate them. That is electrolysis, explained in Chapters 33 and 36. Aqueous electrolysis failed, yielding hydrogen and oxygen because water decomposed first. Davy changed strategy, melting potash to eliminate water's interference, then applying current. On October 6, 1807, silvery droplets appeared at the cathode, some immediately igniting with a pop. He recorded his delight, reportedly jumping about the laboratory and tearing his notebook in excitement. He named the element potassium. Days later, the same method separated soda to produce sodium.

Remember the method, not Davy's luck: \textbf{whether a substance is an element depends not on its stubbornness but on whether you possess a stronger tool for breaking it apart}. New tools can rewrite lists. Conversely, the new metals' ignition in air and explosive water reactions explain their concealment. In nature, they never wait around as elemental metals; they surrender electrons, become quiet ions, and lock themselves into ores and salt lakes. Every grain of salt and piece of rock shows their presence, yet no piece of sodium metal on Earth is natural.
\end{zhengju}

\section{When This Model Fails}

Reading charge from group and increasing activity downward form a useful map, but maps are not territories. At least three boundaries matter.

First, this describes \textbf{the electron-loss tendency of elemental atoms} and fails completely for already ionized elements. \chem{Na^+} can sit quietly in water for millennia because its loss is complete and accounts settled. Active sodium always means sodium metal, not sodium ions in salt.

Second, actual reactions depend on those three conditions. Aluminum ranks ahead of iron in activity, yet aluminum pots hold soup daily while iron nails rust. An oxide film can reverse paper rankings in practice. Before predicting reaction, ask about the surface shell, temperature, and contact area.

Third, sodium's water explosion is subtler than expected. Textbooks describe reaction heat igniting hydrogen, generally adequate for small pieces. High-speed photography, however, reveals countless surface spikes at the first instant, like a tiny blue hedgehog. Some research attributes this to electrons entering water almost instantly, leaving strongly positive metal whose self-repulsion bursts it apart, a Coulomb explosion potentially preceding hydrogen ignition. Mechanistic details remain debated. Science has not written the final page even on this century-old reaction.

Fourth, families contain individuals. Beryllium tops Group 2, so should it be most reactive? Quite the opposite: its dense, robust oxide film tightly encloses it, making the metal quite stable in air. Its compounds, however, are highly toxic, and inhaled beryllium dust severely damages lungs. Barium is contrastingly active, with many soluble compounds; soluble barium salts are extremely poisonous. A group maps personality tendencies, not every member's biography. Particular behavior depends on form, environment, and dose.

\begin{wuqu}
``Taking calcium tablets means eating calcium metal, supplementing that reactive metal.'' This representative mistake overlooks that tablets contain calcium carbonate, citrate, or similar compounds, with no metallic calcium, only \textbf{calcium ions} \chem{Ca^{2+}}. The atoms already surrendered two electrons, settled the account, and lost their temper. Swallowed calcium metal would generate hydrogen and strong base in saliva and stomach fluid, burning the digestive tract: poisoning, not supplementation. Element, elemental substance, and ion are different concepts (Chapter 2 warned of this). The body's calcium requirement means ions; calcium bubbling in water means metal. The medicine-bottle resident and the kerosene-bottle resident share identity but have long parted in personality.
\end{wuqu}

\section{Back to the Calcium Tablet}

We can now answer the opening question.

Why is a tablet mild? Its calcium is \textbf{calcium ions that have already surrendered their electrons}, \chem{Ca^{2+}}, electrostatically bound to carbonate \chem{CO_3^{2-}} in orderly calcium-carbonate crystals (Chapters 18 and 24 explain ionic construction). Electron accounts are settled and the energy slope descended; nothing drives it to fuss with water or air. Unsurrendered elemental metal is violent; already-surrendered ions are mild. One element, two lives.

The ions you swallow do not stay free forever. Most are sent to the bone calcium mine as building material, some remain in blood for signaling, and excess leaves in urine to be replaced by the next meal, maintaining precise balance. The opening mystery is solved: medicine and kerosene bottles contain not two different elements but two electron states of one identity.

This also explains the absence of natural sodium and calcium metals. Their restless elemental forms immediately surrender electrons to oxygen, chlorine, or water, becoming ions in rock salt, limestone, and salt lakes. Almost all Group 1 and 2 elements we encounter are ionic. To see elemental metal, force electrons back with electricity as Davy did, then quickly place it under kerosene and watch it carefully.

\begin{shiyan}
\textbf{A Group 2 Element in Eggshell} (Safe at home.)

Materials: an eggshell or small piece of chalk, both calcium carbonate; white vinegar, about \ud{5}{\%} acetic acid; and a transparent glass.

Procedure: Break the shell into two or three pieces, put them in the glass, and cover with vinegar. Observe immediately, after ten minutes, and after several hours.

What to observe: (1) What appears on the shell at once? Where do bubbles begin and where do they go? (2) After hours, is the shell softer or thinner? (3) Leave overnight: what remains next day?

You see calcium ions changing partners. Acetic acid removes carbonate and releases carbon dioxide, the true identity of the bubbles, while calcium ions dissolve in vinegar. Weak acid demolishes the hard shell. Safety: even weak acid should be kept out of eyes; pour the used liquid down the drain with running water. Think about stomach acid, stronger than acetic acid: it explains dissolved, absorbed calcium tablets and why shells and coral suffer when seawater acidifies (Chapter 34 returns to this).
\end{shiyan}

\section{After Surrendering Electrons, They Support Your Life}

Once Group 1 and 2 elements become ions, dangerous characters turn into ubiquitous backstage workers. Consider four.

\textbf{Bones are a whole calcium mine.} An adult contains about \ud{1}{kg} of calcium. At \ud{40}{g} per mole, that is about 25 moles or $1.5\times 10^{25}$ atoms. Roughly 99\% forms calcium phosphate salts in bones and teeth, serving as structure and reserve. The remaining percent dissolves in blood and cellular fluid yet governs muscle contraction and nerve signaling. Even tiny blood-calcium changes are tightly regulated (Chapter 62 revisits hormonal control). Supplementation is therefore not simply more calcium: dose, absorption, and individual differences matter. Follow doctors, not advertisements, on whether and how much to supplement.

\textbf{A checkup involving drinking barium.} Gastrointestinal X-ray examinations may require a white barium meal, suspended \chem{BaSO_4}. But are barium ions not poisonous? Solubility makes the difference. Almost insoluble barium sulfate passes through the gut, producing almost no free ions while strongly absorbing X-rays to outline digestion for the doctor. One element, soluble salts as poison and poorly soluble salts as medicine. Toxicity and treatment always depend on bodily form and dose.

\textbf{Potassium feeds plants.} Crops need enough potassium that it joins nitrogen and phosphorus as a major fertilizer nutrient. Potassium ions regulate cellular water movement and enzyme activity. Potassium in bananas and potatoes travels from soil to table this way. Plant ash containing potassium carbonate is among the oldest fertilizers; kalium itself derives from plant ash.

\textbf{The lightest metal carries the heaviest electrical load.} Lithium is the lightest metal and among the readiest electron donors. Per gram, it supplies many electrons readily, making it a rechargeable-battery star (Chapter 36 explains). Phone and electric-car batteries put its electron-donating personality into a controlled cage for gradual use. Lithium salts are also longstanding bipolar-disorder medicines with a narrow dosing window requiring medical monitoring, another example of form and dose deciding everything.

Do not forget a hidden thread: sodium's and potassium's surrendered electrons need recipients. In water, water molecules receive them and hydrogen results. What if sodium meets a \textbf{particularly eager electron-taker}? That gives another, more complete and more common combination, table salt. The far-right family of electron-seizing elements opposes this chapter's ready donors. Next chapter visits them.

\begin{tiaozhan}
Reactivity \textbf{can be measured}; we need not compare dramatic videos. Place a metal in a solution of its own ions, connect a circuit, and measure its electron-pushing strength, electrode potential (Chapter 36). Typical values in volts, more negative meaning readier donation: potassium about $-2.93$, calcium $-2.87$, sodium $-2.71$, magnesium $-2.37$, lithium $-3.04$. Two counterintuitive points follow. First, lithium is the strongest donor here yet reacts most gently with water. Greatest tendency does not imply greatest speed: its high melting point, $180.5\,^{\circ}\mathrm{C}$, prevents reaction heat from making a darting bead, leaving small contact area and lower rate. Second, calcium's potential is more negative than sodium's yet its water reaction is milder, partly because sparingly soluble calcium hydroxide coats it like a pot lid. Use instruments instead of spectacle and mechanisms to explain surprises: the chemist's standard approach to rankings. Skipping this does not interrupt the main story.
\end{tiaozhan}

\section*{Try It Yourself}

\begin{sikaoti}
\item Using the ionization-energy price cliff, predict without references which ionization of Group 13 aluminum, with 3 outer electrons, jumps sharply. What is its common ion's charge? Draw your reasoning.
\item Predict whether calcium reacts more vigorously with cold water or dilute hydrochloric acid, considering the strength of its electron-taking partner. What changes if you replace an equal-sized piece with powder?
\item A student says active sodium means sprinkling salt in water must produce bubbles. Identify at least two confused concepts and write a correction.
\item Plot energy against reaction progress for sodium metal plus water and sodium ions plus chloride ions, salt dissolved in water. Show relative starting and ending heights and explain which state sits more comfortably and why.
\item Barium meals use nearly insoluble sulfate rather than soluble chloride. Explain their different safety using ionic form and dose. If soluble barium were mistakenly used, what emergency-treatment direction could chemistry suggest? Discuss principles only.
\item A battery advertisement says lithium's lightness gives large capacity. That is half the explanation. Complete it with electron-donation language: besides lightness, what key personality makes lithium suitable?
\end{sikaoti}
